\documentclass{article}
\usepackage{graphicx} % Required for inserting images
\usepackage{amsmath}
\usepackage{algorithmic}
\usepackage{algorithm}
\usepackage[polish]{babel}
\usepackage[utf8]{inputenc}
\usepackage[T1]{fontenc}
\usepackage{float}
\usepackage{listings}
\usepackage[table]{xcolor}
\usepackage{diagbox}
\usepackage{geometry}

\title{Algorytmy Równoległe - Laboratorium 6}
\author{Adam Staniszewski}

\lstset{
    language=python,
    basicstyle=\ttfamily,
    keywordstyle=\color{blue},
    commentstyle=\color{gray},
    stringstyle=\color{orange},
    backgroundcolor=\color{lightgray!20},
    frame=single,
    captionpos=b
}

\begin{document}

\setlength{\parindent}{0pt}

\maketitle

\section{Zasada działania algorytmów typu branch and bound}
\textbf{Branch and bound} to metoda rozwiązywania problemów kombinatorycznych, która opiera się na przeszukiwawniu przestrzeni rozwiązań przy eliminacji części przestrzeni, które nie mogą zawierać lepszego rozwiązania.

\vspace{1em}
Kluczowe elementy algorytmu:
\begin{itemize}
    \item \textbf{Drzewo przeszukiwań}
    \begin{itemize}
        \item Każdy węzeł drzewa reprezentuje częściowe rozwiązanie (podzbiór miast, które komiwojażer odwiedził, oraz aktualny koszt podróży),
        \item Korzeń drzewa odpowiada sytuacji początkowej (bez odwiedzonych miast, z wyjątkiem miasta startowego).
    \end{itemize}
    \item \textbf{Branch}
    \begin{itemize}
        \item Z każdego węzła generowane są nowe węzły, które reprezentują kolejne możliwe decyzje (odwiedzenie następnych miast),
        \item Dla każdego nowego węzła aktualizowany jest koszt trasy.
    \end{itemize}
    \item \textbf{Bound}
    \begin{itemize}
        \item Algorytm oblicza dolne ograniczenie kosztu dla każdej gałęzi,
        \item Dolne ograniczenie jest oszacowaniem minimalnego możliwego kosztu kontynuowania trasy z danego węzła,
        \item Jeżeli dolne ograniczenie jest wyższe niż koszt najlepszego znalezionego dotąd rozwiązania, gałąź jest odrzucana.
    \end{itemize}
    \item \textbf{Pruning}
    \begin{itemize}
        Jeśli dolne ograniczenie dla danego węzła jest większe niż koszt obecnego najlepszego rozwiązania, algorytm nie przeszukuje dalej tej gałęzi,
    \end{itemize}
\end{itemize}
Algorytm kontynuuje rozgałęzianie i ograniczanie, aż wszystkie gałęzie zostaną albo rozwinięte, albo przycięte.

\section{Kod}
\subsection{Stałe}
\begin{lstlisting}
MAXSIZE = float("inf")
N = 12
DEPTH = 2
AGGLOMERATION_RADIUS = 10
NOISE_LEVEL = 100
GENERATE_INSIDE = True
\end{lstlisting}

\subsection{Generowanie miast}
\begin{lstlisting}
def generate_cities(
        n: int, 
        radius: float = 10, 
        noise_level: float = 42, 
        generate_inside: bool = False
    ) -> list[tuple[float, float]]:
    angles = np.linspace(0, 2 * np.pi, n, endpoint=False)
    angles += np.random.uniform(-noise_level, noise_level, size=n)
    
    if generate_inside:
        radii = np.random.uniform(0, radius, size=n)
    else:
        radii = np.full(n, radius)

    return [(radii[i] * np.cos(angles[i]), 
    radii[i] * np.sin(angles[i])) for i in range(n)]
\end{lstlisting}

\subsection{Generowanie macierzy odległości}
\begin{lstlisting}
def distance_matrix(cities: list[tuple[float, float]]):
    N = len(cities)
    dist_mat: np.ndarray = np.zeros((N, N))

    for i in range(N):
        for j in range(N):
            dist_mat[i][j] = math.sqrt(
                (cities[i][0] - cities[j][0])**2 + 
                (cities[i][1] - cities[j][1])**2
                )

    return dist_mat
\end{lstlisting}

\subsection{Kopiowanie ścieżek}
\begin{lstlisting}
def copy_to_final(curr_path, final_path, N):
    final_path[:N + 1] = curr_path[:]
    final_path[N] = curr_path[0]
\end{lstlisting}

\subsection{Pierwsze minimum}
\begin{lstlisting}
def first_min(adj: list[tuple[float, float]], i: int) -> float:
    min = MAXSIZE
    for j in range(len(adj)):
        if adj[i][j] < min and i != j:
            min = adj[i][j]
    return min
\end{lstlisting}

\subsection{Drugie minimum}
\begin{lstlisting}
def second_min(adj: list[tuple[float, float]], i: int) -> float:
    first = second = MAXSIZE
    for j in range(len(adj)):
        if i == j:
            continue
        if adj[i][j] <= first:
            second = first
            first = adj[i][j]
        elif adj[i][j] <= second and adj[i][j] != first:
            second = adj[i][j]
    return second
\end{lstlisting}

\subsection{Rekursywne TSP}
\begin{lstlisting}
def tsp_recursive(
        adj, 
        curr_bound, 
        curr_weight, 
        level, 
        curr_path, 
        visited, 
        final_path, 
        N, 
        final_res
        ):
    if level == N:
        if adj[curr_path[level - 1]][curr_path[0]] != 0:
            curr_res = curr_weight + adj[curr_path[level - 1]][curr_path[0]]
            if curr_res < final_res[0]:
                copy_to_final(curr_path, final_path, N)
                final_res[0] = curr_res
        return
    
    for i in range(N):
        if adj[curr_path[level-1]][i] != 0 and not visited[i]:
            curr_weight += adj[curr_path[level - 1]][i]
            if level == 1:
                curr_bound -= ((first_min(adj, curr_path[level - 1]) +
                                first_min(adj, i)) / 2)
            else:
                curr_bound -= ((second_min(adj, curr_path[level - 1]) +
                                first_min(adj, i)) / 2)
                
            if curr_bound + curr_weight < final_res[0]:
                visited[i] = True
                curr_path[level] = i
                tsp_recursive(
                    adj, 
                    curr_bound, 
                    curr_weight, 
                    level + 1, 
                    curr_path, 
                    visited, 
                    final_path, 
                    N, 
                    final_res
                )
                visited[i] = False

            curr_weight -= adj[curr_path[level - 1]][i]
            curr_bound += ((second_min(adj, curr_path[level - 1]) +
                            first_min(adj, i)) / 2)
\end{lstlisting}

\subsection{Głowna metoda wywołująca branch and bound}
\begin{lstlisting}
def branch_and_bound_tsp(dist_mat, N):
    final_res = [MAXSIZE]
    curr_bound = 0
    curr_path = [-1] * (N + 1)
    visited = [False] * N

    for i in range(N):
        curr_bound += (first_min(dist_mat, i) + second_min(dist_mat, i))
    curr_bound = math.ceil(curr_bound / 2)

    curr_path[0] = 0
    visited[0] = True

    final_path = [-1] * (N + 1)
    tsp_recursive(
        dist_mat, 
        curr_bound, 
        0, 
        1, 
        curr_path, 
        visited, 
        final_path, 
        N, 
        final_res
    )

    return final_res[0], final_path
\end{lstlisting}

\section{Implementacja równoległa}
\begin{lstlisting}
def branch_and_bound_tsp_parallel(
        cities: list[tuple[float, float]], 
        N: int, 
        depth: int
        ):
    dist_mat = distance_matrix(cities)
    initial_paths = [[0, i] for i in range(1, min(N, depth+1))]

    def process_path(path: list[int]) -> tuple[float, list[int]]:
        visited = [False] * N
        visited[0] = True
        curr_bound = 0
        curr_path = [-1] * (N + 1)
        curr_path[0] = path[0]
        for i in range(N):
            curr_bound += (first_min(dist_mat, i) + second_min(dist_mat, i))
        curr_bound = math.ceil(curr_bound / 2)

        final_res = [MAXSIZE]
        final_path = [-1] * (N + 1)
        tsp_recursive(
            dist_mat, 
            curr_bound, 
            0, 
            1, 
            curr_path, 
            visited, 
            final_path, 
            N, 
            final_res
        )
        return final_res[0], final_path

    comm = MPI.COMM_WORLD
    rank = comm.Get_rank()
    size = comm.Get_size()

    if rank == 0:
        initial_paths = comm.bcast(initial_paths, root=0)

        results = []
        for i, path in enumerate(initial_paths):
            result = process_path(path)
            results.append(result)

        final_res = [MAXSIZE]
        final_path = [-1] * (N + 1)
        for res in results:
            if res[0] < final_res[0]:
                final_res[0] = res[0]
                final_path = res[1]

        return final_res[0], final_path

    return None, None
\end{lstlisting}

\section{Eksperymenty}
W ramach eksperymentów rozważymy problemu rozmiaru $N = 10$ i $N = 15$. Sprawdzona będzie również generacja na i wewnątrz okręgu. Parametry odpowiedzialne za losowość rozłożenia miast (szum) będą stałe dla wszystkich eksperymentów.

\subsection{Implementacja sekwencyjna}

\subsubsection{Rozmiar N = 10}

\paragraph{Miasta na okregu}:
\begin{figure}[H]
    \centering
    \includegraphics[width=0.8\linewidth]{seq_10_10_cities.png}
    \caption{Rozkład miast}
    \label{fig:enter-label}
\end{figure}
\begin{figure}[H]
    \centering
    \includegraphics[width=0.8\linewidth]{seq_10_10_heat.png}
    \caption{Heatmapa odległości między miastami}
    \label{fig:enter-label}
\end{figure}
\begin{figure}[H]
    \centering
    \includegraphics[width=0.8\linewidth]{seq_10_10_route.png}
    \caption{Optymalna ścieżka}
    \label{fig:enter-label}
\end{figure}
Algorytm wybrał ścieżkę:
\[
0 \rightarrow 3 \rightarrow 4 \rightarrow 8 \rightarrow 9 \rightarrow 6 \rightarrow 1 \rightarrow 5 \rightarrow 2 \rightarrow 7 \rightarrow 0
\]
jej całkowity koszt to \textbf{58.8202475232184}. Algorytm wykonał się w \textbf{0.07} sekundy.

\paragraph{Miasta wewnątrz okręgu} :
\begin{figure}[H]
    \centering
    \includegraphics[width=0.8\linewidth]{seq_10_10_cities_inner.png}
    \caption{Rozkład miast wewnątrz okręgu}
    \label{fig:enter-label}
\end{figure}
\begin{figure}[H]
    \centering
    \includegraphics[width=0.8\linewidth]{seq_10_10_heat_inner.png}
    \caption{Heatmapa odległości miast.}
    \label{fig:enter-label}
\end{figure}
\begin{figure}[H]
    \centering
    \includegraphics[width=0.8\linewidth]{seq_10_10_route_inner.png}
    \caption{Optymalna ścieżka}
    \label{fig:enter-label}
\end{figure}
Algorytm wybrał ścieżkę:
\[
0 \rightarrow 2 \rightarrow 8 \rightarrow 3 \rightarrow 7 \rightarrow 6 \rightarrow 1 \rightarrow 9 \rightarrow 5 \rightarrow 4 \rightarrow 0
\]
jej całkowity koszt to \textbf{45.94564798204372}. Algorytm wykonał się w \textbf{0.28} sekundy.

\subsubsection{Rozmiar N = 15}
\paragraph{Miasta na okręgu}
\begin{figure}[H]
    \centering
    \includegraphics[width=0.8\linewidth]{seq_15_15_cities.png}
    \caption{Rozkład miast na okręgu}
    \label{fig:enter-label}
\end{figure}
\begin{figure}[H]
    \centering
    \includegraphics[width=0.8\linewidth]{set_15_15_route.png}
    \caption{Optymalna ścieżka}
    \label{fig:enter-label}
\end{figure}
Algorytm wybrał ścieżkę:
\[
0 \rightarrow 3 \rightarrow 13 \rightarrow 10 \rightarrow 2 \rightarrow 4 \rightarrow 5 \rightarrow 1 \rightarrow 14 \rightarrow 7 \rightarrow 9 \rightarrow 8 \rightarrow 6 \rightarrow 12 \rightarrow 11 \rightarrow 0
\]
jej całkowity koszt to \textbf{45.94564798204372}. Algorytm wykonał się w \textbf{9.43} sekundy.

\paragraph{Miasta wewnątrz okręgu}
\begin{figure}[H]
    \centering
    \includegraphics[width=0.8\linewidth]{seq_15_15_route_inner.png}
    \caption{Rozkład miast wewnątrz okręgu}
    \label{fig:enter-label}
\end{figure}
\begin{figure}[H]
    \centering
    \includegraphics[width=1\linewidth]{set_15_15_route_inner.png}
    \caption{Optymalna ścieżk}
    \label{fig:enter-label}
\end{figure}
Algorytm wybrał ścieżkę:
\[
0 \rightarrow 2 \rightarrow 13 \rightarrow 10 \rightarrow 9 \rightarrow 12 \rightarrow 8 \rightarrow 3 \rightarrow 11 \rightarrow 1 \rightarrow 6 \rightarrow 14 \rightarrow 4 \rightarrow 7 \rightarrow 5 \rightarrow 0
\]
jej całkowity koszt to \textbf{56.89942607195136}. Algorytm wykonał się w \textbf{94.98} sekundy.

\paragraph{Wnioski}:
Widzimy, że rozmiar problemu i rozłożenie miast mają bardzo dużyw wpływ na szybkość wykonania algorytmu.

\subsection{Implementacja równoległa}
Eksperymenty dla wersji równoległej przeprowadzimy dla problemu o rozmiarze $N=15$ oraz dla miast wygenerowanych \textbf{wewnątrz} okręgu przy zmiennym parametrze \textbf{d} i liczbie procesów.

\paragraph{d = 2, 2 procesy}
Execution time: 321.55s \\
Minimum cost: 61.06028603428557 \\
Path taken: [0, 4, 1, 7, 3, 10, 6, 11, 14, 9, 8, 13, 5, 12, 2, 0] \\
\begin{figure}[H]
    \centering
    \includegraphics[width=0.8\linewidth]{par22.png}
    \label{fig:enter-label}
\end{figure}

\paragraph{d = 3, 2 procesy}
Execution time: 20.03s \\
Minimum cost: 48.60273717808537 \\
Path taken: [0, 9, 6, 8, 3, 7, 10, 14, 1, 4, 2, 11, 13, 5, 12, 0] \\
\begin{figure}[H]
    \centering
    \includegraphics[width=0.8\linewidth]{par32.png}
    \label{fig:enter-label}
\end{figure}

\paragraph{d = 4, 2 procesy}
Execution time: 41.67
Minimum cost: 52.608882472531604
Path taken: [0, 1, 10, 5, 6, 12, 9, 2, 8, 11, 4, 7, 14, 13, 3, 0]
\begin{figure}[H]
    \centering
    \includegraphics[width=0.8\linewidth]{par42.png}
    \caption{Enter Caption}
    \label{fig:enter-label}
\end{figure}

\paragraph{d = 2, 4 procesy}
Execution time: 166.67
Minimum cost: 57.183862845001066
Path taken: [0, 3, 5, 7, 10, 4, 13, 9, 14, 6, 11, 1, 8, 12, 2, 0]
\begin{figure}[H]
    \centering
    \includegraphics[width=0.8\linewidth]{par24.png}
    \label{fig:enter-label}
\end{figure}

\paragraph{d = 3, 4 procesy}
Execution time: 42.71
Minimum cost: 44.04535449414117
Path taken: [0, 3, 9, 4, 5, 2, 7, 10, 6, 13, 8, 1, 14, 11, 12, 0]
\begin{figure}[H]
    \centering
    \includegraphics[width=0.8\linewidth]{par34.png}
    \label{fig:enter-label}
\end{figure}

\paragraph{d = 4, 4 procesy}
Execution time: 89.96
Minimum cost: 60.22844999730788
Path taken: [0, 1, 13, 7, 5, 2, 3, 4, 12, 9, 14, 10, 6, 8, 11, 0]
\begin{figure}[H]
    \centering
    \includegraphics[width=0.8\linewidth]{par_44.png}
    \label{fig:enter-label}
\end{figure}

\subsection{Przyspieszenie, efektywność, metryka Karpa-Flatta}
\begin{figure}[H]
    \centering
    \includegraphics[width=0.8\linewidth]{speedup.png}
    \label{fig:enter-label}
\end{figure}
\begin{figure}[H]
    \centering
    \includegraphics[width=0.8\linewidth]{effectivness.png}
    \label{fig:enter-label}
\end{figure}
\begin{figure}[H]
    \centering
    \includegraphics[width=0.8\linewidth]{karp-flatt.png}
    \label{fig:enter-label}
\end{figure}

\end{document}
